\documentclass[10pt,a4paper]{amsart}
\usepackage[margin=19mm]{geometry}
\usepackage{amsmath,amssymb,mathtools}
\usepackage[scr=rsfs,cal=euler]{mathalfa}
\usepackage{xfrac}
\usepackage[unicode,hidelinks,bookmarks=false]{hyperref}
\newcommand{\ie}{i.e.\ }
\newcommand{\cf}{cf.\ }
\newcommand{\resp}{respectively\ }
\newcommand{\Subsection}{\subsection}
\providecommand{\nobreakdash}{\nobreak\hbox{-}}
\setlength{\emergencystretch}{2em}
\multlinegap0pt
\usepackage{amssymb}
\usepackage{float}
\usepackage{enumerate}
\usepackage{xcolor}
\usepackage{tikz}
\newtheorem{theorem}{Theorem}
\newtheorem{prop}{Proposition}
\newcommand\BBN{{\mathbb N}}
\newcommand\CA{{\mathscr A}}
\newcommand\CB{{\mathscr B}}
\newcommand\CH{{\mathscr H}}
\newcommand\lc{{\operatorname{lc}}}
\newcommand\rk{\operatorname{rk}}
\title{Hyperpolygonal arrangements}
\author{Lorenzo Giordani, Paul M\"ucksch, Gerhard R\"ohrle, Johannes Schmitt}
\date{Source: arXiv:2502.02274v4 (CC BY 4.0)}
\begin{document}
\maketitle

\noindent\emph{Source and licence.} Hyperpolygonal arrangements. Version arXiv:2502.02274v4 (CC BY 4.0), distributed by arXiv under the Creative Commons Attribution 4.0 International licence, http://creativecommons.org/licenses/by/4.0/, as recorded in the arXiv licence metadata for this exact version and shown on the abstract page https://arxiv.org/abs/2502.02274v4. Full licence notice: \texttt{licenses/arxiv-2502.02274-CC-BY-4.0.txt}.\par\smallskip

\noindent\textit{Editorial scope.} Counted unit: the article's theorem thm:HAformal (source0.tex lines 372--375). The article's complete original proof of it is delivered with it (source0.tex lines 1000--1009, one \texttt{proof} environment of 10 lines, which the article places away from the statement). No mathematics is written, repaired or re-derived: the statement and the proof are the article's own lines, in the article's own notation, and no retained line is substituted or re-flowed.  Retained uncounted as the source-local context the proof needs: lines 993--996.  Nothing was omitted from any retained span, so the package carries no omission record.  No retained line contains a citation, so the wrapper carries no bibliography and no bibliography entry.  No retained line contains a figure environment, an image-inclusion command or drawing code, so no image is delivered and the package declares no asset dependency.  Editorial formatting only, in addition: an AMS \texttt{amsart} preamble that restates the article's own declarations and macros used by the retained lines, namely the environments \texttt{theorem}, \texttt{prop} on separate counters, together with the standard packages those lines need.  The retained environments print in this document as Theorem 1, Proposition 1 in order of appearance, whereas the article prints the same items with the numbering of its own full document.  The complete native source of the article as distributed in the arXiv e-print for this exact version is bundled unchanged as \texttt{sources/172720/source0.tex}.
\begin{prop}
	\label{prop:lcbasis}
	Let $\CA$ be an arrangement of rank $r$. Suppose $\CB \subseteq \CA$ consists of $r$ hyperplanes such that $r(\CB)=r$ and $\lc(\CB)=\CA$. Then $\CA$ is combinatorially formal.
\end{prop}

\begin{theorem}
	\label{thm:HAformal}
		For any $n \in \BBN$,  $\CH_n$ is combinatorially formal.
\end{theorem}

\begin{proof}[{Proof of Theorem \ref{thm:HAformal}}]
Let $\CB = \{\ker x_1, \ldots, \ker x_{n-1}, \ker(x_1 + \ldots + x_n)\} \subseteq \CH_n$.
Then it is easy to see that successively all $H_I$ for $\varnothing \ne I \subseteq [n]$ belong to the line-closure $\lc(\CB)$ of $\CB$, as follows. 
For $i \in [n]$, set $H_i := \ker x_i$, and 
for $\alpha := x_1 + \ldots + x_n$ set $H_\alpha := \ker \alpha$. 
Then for $I = [n]\setminus \{i\}$, we have
$H_I = \ker \left(\alpha -2 x_i\right) \supset H_\alpha \cap H_i$. Thus for each $I$ of cardinality $n-1$, $H_I$ belongs to $\lc(\CB)$. Now iterate the argument above for $I$ of successively smaller cardinality.
Consequently, $\lc(\CB) =  \CH_n$. Since
	$\rk(\CH_n) = n = |\CB| = \rk(\CB)$, it follows from Proposition \ref{prop:lcbasis}	that $\CH_n$ is combinatorially formal.
\end{proof}

\end{document}
